\uuid{0unE}
\titre{Un perceptron à trois entrées : vote à la majorité}
\niveau{L1}
\module{OPTRN}
\chapitre{Réseaux de neurones}
\sousChapitre{Perceptron}
\theme{Perceptron affine, fonctions logiques, inférence}
\auteur{Maxime Nguyen}
\datecreate{2026-10-20}
\organisation{AMSCC}
\difficulte{1}
\contenu{
\texte{Trois capteurs renvoient chacun $0$ (pas d'alerte) ou $1$ (alerte). Un perceptron à trois entrées $x_1,x_2,x_3\in\{0,1\}$ est défini par les poids $1$, $1$ et $1$ et par le biais $-2$. Sa fonction d'activation est la fonction de Heaviside $H$, avec $H(t)=1$ si $t\geq0$ et $H(t)=0$ sinon :
\[
F(x_1,x_2,x_3)=H(x_1+x_2+x_3-2).
\]}
\begin{enumerate}
\item \question{Calculer $F(1,0,1)$, $F(0,0,1)$ et $F(1,1,1)$.}
\indication{Calculer d'abord la somme $x_1+x_2+x_3-2$.}
\reponse{On a $F(1,0,1)=H(0)=1$, $F(0,0,1)=H(-1)=0$ et $F(1,1,1)=H(1)=1$.}
\item \question{Dresser le tableau des valeurs de $F$ sur les huit entrées de $\{0,1\}^3$.}
\reponse{\begin{center}
\begin{tabular}{ccc|cc}
$x_1$ & $x_2$ & $x_3$ & $x_1+x_2+x_3-2$ & $F$\\\hline
$0$ & $0$ & $0$ & $-2$ & $0$\\
$1$ & $0$ & $0$ & $-1$ & $0$\\
$0$ & $1$ & $0$ & $-1$ & $0$\\
$0$ & $0$ & $1$ & $-1$ & $0$\\
$1$ & $1$ & $0$ & $0$ & $1$\\
$1$ & $0$ & $1$ & $0$ & $1$\\
$0$ & $1$ & $1$ & $0$ & $1$\\
$1$ & $1$ & $1$ & $1$ & $1$
\end{tabular}
\end{center}}
\item \question{Décrire en une phrase ce que calcule ce perceptron.}
\reponse{La sortie vaut $1$ si et seulement si au moins deux des trois capteurs signalent une alerte : c'est un vote à la majorité.}
\item \question{On garde les poids $1$, $1$, $1$ et on change seulement le biais. Quelle fonction logique obtient-on avec le biais $-3$ ? avec le biais $-1$ ?}
\indication{Le perceptron renvoie $1$ si et seulement si $x_1+x_2+x_3\geq k$ pour un certain entier $k$.}
\reponse{Avec le biais $-3$, la sortie vaut $1$ si et seulement si $x_1+x_2+x_3\geq3$, c'est-à-dire si les trois entrées valent $1$ : c'est le « ET ». Avec le biais $-1$, la sortie vaut $1$ si et seulement si $x_1+x_2+x_3\geq1$, c'est-à-dire si au moins une entrée vaut $1$ : c'est le « OU ».}
\end{enumerate}
}
